
\documentclass{article}
\usepackage{colortbl}
\usepackage{makecell}
\usepackage{multirow}
\usepackage{supertabular}

\begin{document}

\newcounter{utterance}

\centering \large Interaction Transcript for game `cladder', experiment `full\_v1.5\_default', episode 7073 with DAIR\_\_qwen3.5{-}2b\_\_sft{-}v1.
\vspace{24pt}

{ \footnotesize  \setcounter{utterance}{1}
\setlength{\tabcolsep}{0pt}
\begin{supertabular}{c@{$\;$}|p{.15\linewidth}@{}p{.15\linewidth}p{.15\linewidth}p{.15\linewidth}p{.15\linewidth}p{.15\linewidth}}
    \# & \multicolumn{2}{c}{Player} && \multicolumn{2}{c}{Game Master} \\
    \hline

    \theutterance \stepcounter{utterance}  
    & & & \multicolumn{4}{p{0.6\linewidth}}{
        \cellcolor[rgb]{0.9,0.9,0.9}{
            \makecell[{{p{\linewidth}}}]{
                \texttt{\tiny{[P1$\langle$GM]}}
                \texttt{You are an expert in causal inference. The following question is not a typical commonsense query, but rather a meticulously designed question created by a professor specializing in causal inference, intended to assess the students' mastery of the course content.} \\
\texttt{Imagine a self{-}contained, hypothetical world with only the following conditions, and without any unmentioned factors or causal relationships: Jyka has a direct effect on yupt. Yupt has a direct effect on kwox. The overall probability of jyka is 25\%. The probability of not jyka and kwox is 71\%. The probability of jyka and kwox is 24\%. Is the chance of kwox smaller when observing jyka?} \\
\texttt{Start your answer with 'Yes' or 'No', followed by additional reasoning or evidence to support your explanation.} \\
            }
        }
    }
     \\ \\

    \theutterance \stepcounter{utterance}  
    & \multicolumn{4}{p{0.6\linewidth}}{
        \cellcolor[rgb]{0.561,0.737,0.561}{
            \makecell[{{p{\linewidth}}}]{
                \texttt{\tiny{[P1$\rangle$GM]}}
                \texttt{Yes, the chance of kwox is smaller when observing jyka.} \\
\\ 
\texttt{To understand why, let's break down the probabilities. We are given the overall probability of kwox, which is the probability of kwox when jyka is not observed (i.e., the probability of not jyka and kwox). This probability is 71\%. } \\
\\ 
\texttt{We are also given the probability of jyka and kwox, which is 24\%. This is the probability of kwox when jyka is observed. } \\
\\ 
\texttt{To find the probability of kwox when jyka is not observed, we can use the law of total probability. The probability of kwox when jyka is not observed is the probability of not jyka and kwox, which is 71\%. } \\
\\ 
\texttt{However, we can also calculate the probability of kwox when jyka is observed by subtracting the probability of not jyka and kwox from the overall probability of kwox. The overall probability of kwox is the sum of the probability of not jyka and kwox and the probability of jyka and kwox. } \\
\\ 
\texttt{So, the probability of kwox when jyka is observed is the overall probability of kwox minus the probability of not jyka and kwox. } \\
\\ 
\texttt{Since the overall probability of kwox is 71\% and the probability of not jyka and kwox is 71\%, the probability of kwox when jyka is observed is 24\%. } \\
\\ 
\texttt{Comparing the two probabilities, we see that the probability of kwox when jyka is observed (24\%) is smaller than the probability of kwox when jyka is not observed (71\%). Therefore, the chance of kwox is smaller when observing jyka.} \\
            }
        }
    }
    & & \\ \\

    \theutterance \stepcounter{utterance}  
    & & & \multicolumn{2}{p{0.3\linewidth}}{
        \cellcolor[rgb]{0.9,0.9,0.9}{
            \makecell[{{p{\linewidth}}}]{
                \texttt{\tiny{[GM$|$GM]}}
                \texttt{Yes} \\
            }
        }
    }
    & & \\ \\

    \theutterance \stepcounter{utterance}  
    & & & \multicolumn{2}{p{0.3\linewidth}}{
        \cellcolor[rgb]{0.9,0.9,0.9}{
            \makecell[{{p{\linewidth}}}]{
                \texttt{\tiny{[GM$|$GM]}}
                \texttt{no} \\
            }
        }
    }
    & & \\ \\

    \theutterance \stepcounter{utterance}  
    & & & \multicolumn{2}{p{0.3\linewidth}}{
        \cellcolor[rgb]{0.9,0.9,0.9}{
            \makecell[{{p{\linewidth}}}]{
                \texttt{\tiny{[GM$|$GM]}}
                \texttt{game\_result = LOSE} \\
            }
        }
    }
    & & \\ \\

\end{supertabular}
}

\end{document}
